\documentclass{article}
\usepackage{graphicx} % Required for inserting images

\title{A Quantitative Framework for Rational HEPHA-RNA Design}
\author{子恒 魏}
\date{October 2026}

\begin{document}

\maketitle




\section{Polymer Conformational Model}

We approximate the mRNA as a freely jointed chain (FJC):
\[
\boldsymbol{\rho}
=
\sum_{i=1}^{N}\mathbf b_i,
\qquad
|\mathbf b_i|=b,
\]
where \(b\) is the Kuhn length and \(N=L/b\) is the number of Kuhn segments.

For sufficiently large \(N\), the end-to-end vector can be approximated by a Gaussian distribution:
\[
\boxed{
p(\boldsymbol{\rho})
=
\left(
\frac{3}{2\pi Nb^2}
\right)^{3/2}
\exp\left(
-\frac{3\rho^2}{2Nb^2}
\right)
}
\]

where \(p(\boldsymbol{\rho})\) is the three-dimensional probability density of the end-to-end vector and
\(\rho=|\boldsymbol{\rho}|\).

Since the protein-mediated bridge depends primarily on the distance between the two interaction sites, we integrate over all orientations to obtain the radial probability density:
\[
\boxed{
p_{\rho}(\rho)
=
4\pi \rho^2 p(\boldsymbol{\rho})
}
\]

or explicitly,

\[
p_{\rho}(\rho)
=
4\pi \rho^2
\left(
\frac{3}{2\pi Nb^2}
\right)^{3/2}
\exp\left(
-\frac{3\rho^2}{2Nb^2}
\right).
\]

We assume that protein-mediated bridging is geometrically possible when the distance between the two interaction sites lies within a capture range:
\[
\rho_{\min}\leq \rho\leq \rho_{\max}.
\]

The equilibrium fraction of RNA conformations within this protein-compatible range is therefore

\[
\boxed{
P_{\rm enc}^{\rm loop}
=
\int_{\rho_{\min}}^{\rho_{\max}}
p_{\rho}(\rho)\,d\rho
}
\]

where \(P_{\rm enc}^{\rm loop}\) is dimensionless.

The corresponding capture region is a spherical shell with volume

\[
\boxed{
V_{\rm loop}
=
\frac{4\pi}{3}
\left(
\rho_{\max}^3-\rho_{\min}^3
\right)
}
\]

We define the average spatial probability density within this capture shell as

\[
\boxed{
\bar p_{\rm loop}
=
\frac{P_{\rm enc}^{\rm loop}}
{V_{\rm loop}}
}
\]

which has units of inverse volume.

To convert this spatial probability density into an effective concentration, we use the relation between the spatial density of one molecule and its molar concentration:

\[
C_{\rm eff}^{\rm loop}
=
\frac{\bar p_{\rm loop}}{N_A},
\]

with the appropriate unit conversion between volume and molarity.

Thus,

\[
\boxed{
C_{\rm eff}^{\rm loop}
=
\frac{1}{N_A}
\frac{
\displaystyle
\int_{\rho_{\min}}^{\rho_{\max}}
4\pi \rho^2p(\boldsymbol{\rho})\,d\rho
}{
\displaystyle
\frac{4\pi}{3}
\left(
\rho_{\max}^3-\rho_{\min}^3
\right)
}
}
\]

where \(C_{\rm eff}^{\rm loop}\) represents the average effective concentration of one interaction site around the other within the protein-compatible capture shell.

\section{Protein-Mediated Bridging}

We next consider protein-mediated bridging between the 5' and 3' ends of the mRNA, for example through interactions involving PABP and eIF4G.

The standard-state bimolecular association constant is

\[
\boxed{
K_{\rm bind}^{\circ}
=
\frac{1}{C^\circ}
\exp\left(
-\frac{\Delta G_{\rm bind}^{\circ}}{RT}
\right)
}
\]

where \(C^\circ=1\,\mathrm{M}\) is the standard concentration.

The intramolecular looping equilibrium is then described by

\[
\boxed{
K_{\rm loop}
=
K_{\rm bind}^{\circ}
C_{\rm eff}^{\rm loop}
}
\]

or equivalently,

\[
\boxed{
K_{\rm loop}
=
\frac{C_{\rm eff}^{\rm loop}}{C^\circ}
\exp\left(
-\frac{\Delta G_{\rm bind}^{\circ}}{RT}
\right)
}
\]

Finally, treating the mRNA as a two-state system,

\[
\mathrm{Open}
\rightleftharpoons
\mathrm{Closed},
\]

the equilibrium fraction of closed-loop conformations is

\[
\boxed{
f_{\rm closed}
=
\frac{K_{\rm loop}}
{1+K_{\rm loop}}
}
\]

\section{Closed-Loop Positional Model}

Once the mRNA adopts the closed-loop state, we approximate the mRNA as a Gaussian ring and consider the spatial encounter between the BD and the 5' translation-initiation machinery.

Let \(n\) denote the number of Kuhn segments along one contour path between the BD and the 5' translation-initiation anchor. The complementary contour path therefore contains \(N-n\) Kuhn segments. The spatial separation between the two sites is denoted by \(R\).

For a Gaussian ring, the three-dimensional probability density of the spatial separation \(R\) is

\[
\boxed{
p_{\rm ring}(\mathbf R\mid n)
=
\left[
\frac{3N}
{2\pi b^2n(N-n)}
\right]^{3/2}
\exp\left[
-\frac{3NR^2}
{2b^2n(N-n)}
\right]
}
\]

where \(R=|\mathbf R|\).

Integrating over all orientations gives the corresponding radial probability density:

\[
\boxed{
p_{\rm ring}(R\mid n)
=
4\pi R^2
\left[
\frac{3N}
{2\pi b^2n(N-n)}
\right]^{3/2}
\exp\left[
-\frac{3NR^2}
{2b^2n(N-n)}
\right]
}
\]

Because the interaction between the BD and the 5' translation-initiation machinery involves finite-sized molecular interfaces rather than an ideal point contact, we assume that productive interaction occurs when their spatial separation lies within a finite capture range,

\[
R_{\min}\leq R\leq R_{\max}.
\]

The probability of finding the BD within this capture range is therefore

\[
\boxed{
P_{\rm enc}^{\rm ring}(n)
=
\int_{R_{\min}}^{R_{\max}}
p_{\rm ring}(R\mid n)\,dR
}
\]

with the corresponding capture-shell volume

\[
\boxed{
V_{\rm BD}
=
\frac{4\pi}{3}
\left(
R_{\max}^3-R_{\min}^3
\right)
}
\]

We define the shell-averaged effective concentration of the BD around the 5' translation-initiation machinery as

\[
\boxed{
C_{\rm eff}^{\rm ring}(n)
=
\frac{1}{N_A}
\frac{
P_{\rm enc}^{\rm ring}(n)
}{
V_{\rm BD}
}
}
\]

with the appropriate unit conversion between volume and molarity.

Equivalently,

\[
\boxed{
C_{\rm eff}^{\rm ring}(n)
=
\frac{1}{N_A V_{\rm BD}}
\int_{R_{\min}}^{R_{\max}}
4\pi R^2
\left[
\frac{3N}
{2\pi b^2n(N-n)}
\right]^{3/2}
\exp\left[
-\frac{3NR^2}
{2b^2n(N-n)}
\right]dR
}
\]

This quantity describes the positional contribution of the polymer conformation to the interaction between the BD and the 5' translation-initiation machinery.

\subsection{cis Configuration}

In the cis configuration, the BD is covalently incorporated into the same mRNA molecule as the 5' translation-initiation machinery. Therefore, no additional intermolecular BD--RNA binding free energy is required.

The translational effect is consequently determined by the probability of forming the closed-loop state and the positional effective concentration of the BD:

\[
\boxed{
TE_{\rm cis}(n)
\propto
f_{\rm closed}
C_{\rm eff}^{\rm ring}(n)
}
\]

Thus, under the assumption that the presence of the BD does not substantially perturb the open--closed equilibrium, the positional dependence of the cis configuration is determined by

\[
\boxed{
TE_{\rm cis}(n)
\propto
C_{\rm eff}^{\rm ring}(n)
}
\]

where \(n\) specifies the position of the BD along the closed-loop mRNA.

\subsection{trans Configuration}

In the trans configuration, the HEPHA-RNA carrying the BD is an independent RNA molecule. In addition to the spatial encounter between the BD and the 5' translation-initiation machinery, productive regulation therefore requires binding of the BD to its target site on the mRNA.

The standard-state association constant for BD--target RNA binding is

\[
\boxed{
K_{\rm BD}^{\circ}(n)
=
\frac{1}{C^\circ}
\exp\left[
-\frac{\Delta G_{\rm BD}^{\circ}(n)}
{RT}
\right]
}
\]

where \(C^\circ=1\,\mathrm{M}\) is the standard concentration and
\(\Delta G_{\rm BD}^{\circ}(n)\) is the standard binding free energy between the BD and the target site at position \(n\).

Importantly, \(\Delta G_{\rm BD}^{\circ}(n)\) is explicitly position-dependent because different target sites along the mRNA may have different sequence and structural properties, leading to different binding affinities.

The spatial encounter and molecular binding contributions can then be combined into an effective intramolecular association term:

\[
\boxed{
K_{\rm trans}(n)
=
K_{\rm BD}^{\circ}(n)
C_{\rm eff}^{\rm ring}(n)
}
\]

or equivalently,

\[
\boxed{
K_{\rm trans}(n)
=
\frac{
C_{\rm eff}^{\rm ring}(n)
}{
C^\circ
}
\exp\left[
-\frac{\Delta G_{\rm BD}^{\circ}(n)}
{RT}
\right]
}
\]

This expression separates the two contributions to trans-acting regulation: \(C_{\rm eff}^{\rm ring}(n)\) describes the positional and polymer-conformational contribution, whereas \(\Delta G_{\rm BD}^{\circ}(n)\) describes the sequence- and structure-dependent binding contribution.

Consequently, the translational effect of a trans-acting BD targeting site \(n\) can be written as

\[
\boxed{
TE_{\rm trans}(n)
\propto
f_{\rm closed}
C_{\rm eff}^{\rm ring}(n)
K_{\rm BD}^{\circ}(n)
}
\]

or

\[
\boxed{
TE_{\rm trans}(n)
\propto
f_{\rm closed}
\frac{
C_{\rm eff}^{\rm ring}(n)
}{
C^\circ
}
\exp\left[
-\frac{\Delta G_{\rm BD}^{\circ}(n)}
{RT}
\right]
}
\]

Thus, the model naturally decomposes the positional effect of a trans-acting element into two independent contributions: the spatial accessibility determined by polymer conformational statistics and the binding affinity determined by the free energy of BD--target RNA interaction.

For comparing different target sites on the same mRNA, \(f_{\rm closed}\) is a common factor and can therefore be omitted. The site-dependent contribution to trans-acting regulation is then

\[
\boxed{
TE_{\rm trans}(n)
\propto
C_{\rm eff}^{\rm ring}(n)
K_{\rm BD}^{\circ}(n)
}
\]

or equivalently,

\[
\boxed{
TE_{\rm trans}(n)
\propto
\frac{
C_{\rm eff}^{\rm ring}(n)
}{
C^\circ
}
\exp\left[
-\frac{\Delta G_{\rm BD}^{\circ}(n)}
{RT}
\right]
}
\]

Thus, when comparing candidate BD binding sites, the model contains two site-dependent contributions:

\[
\boxed{
\text{Site score}(n)
\propto
\underbrace{
C_{\rm eff}^{\rm ring}(n)
}_{\text{polymer positional contribution}}
\;
\underbrace{
\exp\left[
-\frac{\Delta G_{\rm BD}^{\circ}(n)}
{RT}
\right]
}_{\text{binding contribution}}
}
\]

The first term captures the positional accessibility of the target site within the closed-loop mRNA, whereas the second term captures the sequence- and structure-dependent affinity of the BD for that site.

\subsection{Constrained Closed-Loop Model}

The Gaussian-ring model assumes that translational enhancement
depends only on the spatial proximity between the BD and the
translation-initiation machinery.

However, experimental observations suggest that proximity alone
is insufficient to explain the positional dependence of the BD.
Instead, optimal activity appears when the effective length of
the 3'-side arm approximately matches that of the 5'-side arm.

We therefore introduce an additional conformational constraint
that represents the geometric restrictions imposed by the
5'--3' closed-loop architecture.

Let

\[
n_{5A}
\]

denote the effective conformational length between the 5' end
and the AUG codon, and

\[
n_{B3}(n)
\]

denote the effective conformational length between the BD site
and the poly(A) tail.

These quantities do not necessarily correspond to the actual
contour lengths of the RNA, but instead represent coarse-grained
effective lengths arising from RNA structure and protein-mediated
organization.

We assume that the closed-loop architecture preferentially
stabilizes conformations in which the two arms have similar
effective lengths.

This preference is represented by a harmonic conformational
potential:

\[
\boxed{
U_{\rm match}(n)
=
\frac{k_{\rm match}}{2}
\left[
n_{5A}
-
n_{B3}(n)
\right]^2
}
\]

where \(k_{\rm match}\) represents the strength of the
conformational constraint.

According to the Boltzmann distribution, the statistical weight
associated with this constraint is

\[
\boxed{
P_{\rm match}(n)
=
\exp\left[
-\frac{
U_{\rm match}(n)
}{RT}
\right]
}
\]

or explicitly,

\[
\boxed{
P_{\rm match}(n)
=
\exp\left[
-\frac{
k_{\rm match}
}{
2RT
}
\left(
n_{5A}
-
n_{B3}(n)
\right)^2
\right]
}
\]

This factor describes the probability that a BD position is
compatible with the global closed-loop architecture.

The effective concentration in the constrained closed-loop model
is therefore

\[
\boxed{
C_{\rm eff}^{\rm constrained}(n)
=
C_{\rm eff}^{\rm ring}(n)
P_{\rm match}(n)
}
\]

which gives

\[
\boxed{
C_{\rm eff}^{\rm constrained}(n)
=
C_{\rm eff}^{\rm ring}(n)
\exp\left[
-\frac{
k_{\rm match}
}{
2RT
}
\left(
n_{5A}
-
n_{B3}(n)
\right)^2
\right]
}
\]

Finally, the trans-site score becomes

\[
\boxed{
\text{Site score}(n)
\propto
C_{\rm eff}^{\rm constrained}(n)
\exp\left[
-\frac{\Delta G_{\rm BD}^{\circ}(n)}
{RT}
\right]
}
\]

or equivalently,

\[
\boxed{
\text{Site score}(n)
\propto
C_{\rm eff}^{\rm ring}(n)
\exp\left[
-\frac{
k_{\rm match}
}{
2RT
}
\left(
n_{5A}
-
n_{B3}(n)
\right)^2
\right]
\exp\left[
-\frac{\Delta G_{\rm BD}^{\circ}(n)}
{RT}
\right]
}
\]

\subsection{Effector Domain Kinetic Model}

While the BD determines where HEPHA-RNA interacts with the target mRNA, the ED determines how efficiently translational machinery is recruited. We therefore model ED activity as a kinetic balance between ribosome capture and release.

For an ED designed against a ribosomal protein, we use the predicted interface confidence, \(\mathrm{ipTM}\), as a proxy for the strength of the ED--protein interaction. Because \(\mathrm{ipTM}\) is not itself a thermodynamic quantity, we introduce a latent interaction-strength parameter,

\[
\boxed{
\theta_{\rm ED}
=
w_{\rm ED}\,\mathrm{ipTM}
+
b_{\rm ED}
}
\]

where \(w_{\rm ED}\) and \(b_{\rm ED}\) are parameters calibrated from experimental data.

We then consider the reversible binding process between the ED and its ribosomal target \(R\),

\[
E+R
\underset{k_{\rm off}}{\overset{k_{\rm on}}{\rightleftharpoons}}
ER,
\]

where \(E\) denotes the ED and \(ER\) denotes the ED--ribosome complex. The corresponding capture and release rates are

\[
\boxed{
v_{\rm capture}
=
k_{\rm on}[E][R]
}
\]

and

\[
\boxed{
v_{\rm release}
=
k_{\rm off}[ER].
}
\]

To connect the predicted interaction strength with the kinetic parameters, we assume that stronger ED--ribosome interactions primarily reduce the dissociation rate,

\[
\boxed{
k_{\rm off}
=
k_{\rm off}^{0}
\exp\left(
-\beta\theta_{\rm ED}
\right)
}
\]

where \(k_{\rm off}^{0}\) is a reference dissociation rate and \(\beta\) controls the sensitivity of the kinetic model to the predicted interaction strength. Under this formulation, increasing \(\mathrm{ipTM}\) corresponds to a longer ED--ribosome residence time,

\[
\boxed{
\tau_{\rm residence}
=
\frac{1}{k_{\rm off}}.
}
\]

However, productive recruitment requires both efficient capture and subsequent release of the translational machinery. The probability of ribosome capture can be approximated by the occupancy of the ED--ribosome complex,

\[
\boxed{
P_{\rm capture}
=
\frac{
k_{\rm on}[R]
}{
k_{\rm on}[R]+k_{\rm off}
}.
}
\]

The release process is governed by \(k_{\rm off}\). We therefore define the productive recruitment efficiency as

\[
\boxed{
\Phi_{\rm ED}
\propto
P_{\rm capture}\,k_{\rm off}
}
\]

which gives

\[
\boxed{
\Phi_{\rm ED}
\propto
\frac{
k_{\rm on}[R]\,k_{\rm off}
}{
k_{\rm on}[R]+k_{\rm off}
}.
}
\]

This formulation captures the kinetic trade-off inherent to ED design. When the ED--ribosome interaction is too weak, \(k_{\rm off}\) is large and the ribosome cannot be efficiently retained after capture. Conversely, when the interaction is excessively strong, \(k_{\rm off}\) becomes too small, increasing the residence time and limiting ribosome release. Productive recruitment therefore occurs at an intermediate interaction strength rather than at the maximum possible affinity.

Thus, the ED design objective can be summarized as

\[
\boxed{
\mathrm{ipTM}
\longrightarrow
\theta_{\rm ED}
\longrightarrow
(k_{\rm on},k_{\rm off})
\longrightarrow
\Phi_{\rm ED}
}
\]

where \(\Phi_{\rm ED}\) provides a kinetic measure of productive ribosome recruitment rather than simply maximizing binding affinity.


\subsection{Joint Optimization of BD and ED}

The BD and ED contribute to HEPHA-RNA activity through distinct but complementary mechanisms. The BD determines the positional accessibility and target-binding efficiency of the HEPHA-RNA, whereas the ED determines the kinetic efficiency of translational machinery recruitment. We therefore combine the two components into a unified design score.

For a given BD position \(n\) and ED design characterized by its predicted interface confidence \(\mathrm{ipTM}\), the overall score is defined as

\[
\boxed{
S_{\rm HEPHA}(n,\mathrm{ipTM})
\propto
S_{\rm BD}(n)\,
\Phi_{\rm ED}(\mathrm{ipTM})
}
\]

where the BD contribution is

\[
\boxed{
S_{\rm BD}(n)
=
C_{\rm eff}^{\rm constrained}(n)
\exp\left[
-\frac{\Delta G_{\rm BD}^{\circ}(n)}
{RT}
\right]
}
\]

and the ED contribution is

\[
\boxed{
\Phi_{\rm ED}(\mathrm{ipTM})
\propto
\frac{
k_{\rm on}[R]\,k_{\rm off}(\mathrm{ipTM})
}{
k_{\rm on}[R]+k_{\rm off}(\mathrm{ipTM})
}.
}
\]

Thus,

\[
\boxed{
S_{\rm HEPHA}(n,\mathrm{ipTM})
\propto
C_{\rm eff}^{\rm constrained}(n)
\exp\left[
-\frac{\Delta G_{\rm BD}^{\circ}(n)}
{RT}
\right]
\Phi_{\rm ED}(\mathrm{ipTM})
}
\]

describes the combined effect of target-site accessibility, BD binding, and ED-mediated ribosome recruitment.

For a fixed target mRNA, the sequence-dependent BD affinity and the global closed-loop parameters can be treated as fixed when examining the joint dependence on BD position and ED interaction strength. The resulting score can therefore be approximated as

\[
\boxed{
S_{\rm HEPHA}(n,\mathrm{ipTM})
\approx
S_{\rm BD}(n)\,
\Phi_{\rm ED}(\mathrm{ipTM}).
}
\]

This separability produces two independent optima: an optimal BD position,

\[
\boxed{
n^*
=
\arg\max_n S_{\rm BD}(n),
}
\]

and an optimal ED interaction strength,

\[
\boxed{
\mathrm{ipTM}^*
=
\arg\max_{\mathrm{ipTM}}
\Phi_{\rm ED}(\mathrm{ipTM}).
}
\]

The joint optimum is therefore

\[
\boxed{
(n^*,\mathrm{ipTM}^*)
=
\arg\max_{n,\mathrm{ipTM}}
S_{\rm HEPHA}(n,\mathrm{ipTM}).
}
\]

Importantly, neither component is expected to favor the strongest possible interaction indefinitely. The constrained-ring model produces a positional optimum for the BD, while the capture--release kinetics of the ED produces an intermediate optimal interaction strength. Consequently, the joint design landscape is expected to exhibit a localized maximum in the two-dimensional design space spanned by BD position and ED interaction strength.

This framework therefore transforms HEPHA-RNA design from independent optimization of individual domains into a unified, experimentally calibratable fitness landscape.
\end{document}